% Nouveau EGA V (cote 25): archive cote 25, site page 177 (PDF page 178).
% typescript. Draft compared: none.
% AI-transcribed by Claude Opus 5.5 (claude-opus-5-5, adaptive thinking, effort high), Message Batches API, batch ega25b (msgbatch_01A7n5YuXubUPNK1zsepi8RG), 2026-11-13; not yet checked by eye.
\scanpage{25/177}
\mspage{\hand{4}}

(c) Si E est partout de rang impair, on a $2\hand{$\widetilde{\phantom{d}}$}\!d(Q)^{\unsure{\hand{2}}} = d(Q)$.

\xout{\unsure{De plus, le discriminant corrigé d a ainsi introduit satisfait à}}

\xout{\unsure{les conditions suivantes}} \quad NB Il n'est pas raisonnable d'énoncer la

propr\over{\illegible{1}}{\unsure{ié}}té de compatibilité avec le changement de base, s ans énoncer en

même temps\hand{,} ou même avant\hand{,} la propriété de fonctorialité par isomorphismes,

étant donné que le changement de base \inserted{sur $\underline{E}$} lui-même n'est \over{p}{g}ère défini que

modulo isomorpphisme; il n'en s erait plus de même si on se restreignait

au cas $E = \underline{O}_S{}^n$ ,(ce qui ne serait pas commode pour les références).

\xout{Corollaire 16.7.} \quad \uline{Définition} 16.7. La section $\hand{$\widetilde{\phantom{d}}$}\!d(Q)$ de $\text{dét}(E)^{\otimes 2}$

s'appelera le \uline{discriminant corrigé} de la forme quadratique Q.

\uline{Corollaire} 16.8. \unsure{Soinnt|Soient} Q,Q' des formes relatives à $\underline{E}$, $\underline{E}'$ , supposons

la parité des rangs de $\underline{E}$ (resp. $\underline{E}'$) constante sur S. Alors

a) Si $\underline{E}$ et $\underline{E}'$ ne sont pas tous deux de rang impair, on a
\[ \hand{$\widetilde{\phantom{d}}$}\!d(Q \oplus Q') = \hand{$\widetilde{\phantom{d}}$}\!d(Q)\hand{$\widetilde{\phantom{d}}$}\!d(Q') \]

b) Si $\underline{E}$ et $\underline{E}'$ sont tous d eux de rang impair, on a
\[ \hand{$\widetilde{\phantom{d}}$}\!d(Q \oplus Q') = \hand{2}\,\hand{$\widetilde{\phantom{d}}$}\!d(Q)\hand{$\widetilde{\phantom{d}}$}\!d(Q') \]

La vérification est triviale.

Nous pouvons maintenant énoncer le résultat principal du présent

numéro:

\uline{Théorème} 16.9. Soient $\underline{E}$ un Module loc libre de rang fini sur le préschém-%

ma S, $Q \in \Gamma\,\uline{Sym}^2(\underline{E})$. Les conditions suivantes sont équivalentes:

(i) Q est lisse.

(ii) Le d\unsure{\over{s}{i}}scriminant modifié $\hand{$\widetilde{\phantom{d}}$}\!d(Q) \in \Gamma\,\text{dét}(\underline{E})^{\otimes 2}$ est inversible.

(iii) \inserted[typed]{(Lorsque $\underline{E}$ est d e rang constant\xout{s} n)} Il existe un morphisme surjectif S' \quad S tel que la forme $Q_{S'}$

déduite de Q par changement de base soit isomorphe à la forme quadrati-%

que standart en n variables.

(i\xout{x}i bis). Comme (iii), mais en supposant S'\hand{$\rightarrow$}S fidèlement plat de

présentation finie.

\mspage{\hand{177}}
